% Source: https://artofproblemsolving.com/community/c3h2024348_number_bases
% Stable identity: aops-topic-2024348-q01
Each base-$4$ digit corresponds to two binary digits: $2_4=10_2$ and $3_4=11_2$. Since all four base-$4$ digits are $2$ or $3$, the binary form of $n$ is obtained by writing these pairs side by side; it has exactly $8$ digits (the leading pair starts with $1$), and its number of $1$s is $4$ plus the number of base-$4$ digits equal to $3$.

So exactly two of the four base-$4$ digits are $3$ and the other two are $2$. The largest such number puts the $3$s first:
\[
n=3322_4=3\cdot64+3\cdot16+2\cdot4+2=\boxed{250}.
\]
(Check: $250=11111010_2$, which has $8$ digits and six $1$s.)
